% Work/MAT — 1. ročník — hodina 59: Co je funkce?
% Návaznost: 58_ALG_OpakovaniAOverovaniUprav; 60_FNC_DefinicniOborAOborHodnot.
% Autorova učebnice: ucebnice/02_Funkce_a_rovnice/Zaklady_funkci/Funkce.tex,
%   zejména „Co je funkce“ a „Způsoby zadání funkce“.
% Metodické zdroje: Realisticky.cz 020103 Zobrazení,
%   020104 Funkce, definiční obor, 020105 Graf funkce I.
% Formulace, diagramy, graf a úlohy jsou vlastní; externí lekce nejsou opisovány.
\documentclass[11pt,a4paper]{article}
% Spolecny styl pro ucitelske pripravy MAT (Work/MAT)
% Verze 1.1, 2026-10-23
\NeedsTeXFormat{LaTeX2e}
\ProvidesPackage{priprava}[2026/10/23 v1.1 Sablona ucitelskych priprav MAT]

% Zakladni typografie a cestina
\RequirePackage[T1]{fontenc}
\RequirePackage[utf8]{inputenc}
\RequirePackage[czech]{babel}
\RequirePackage{lmodern}
\RequirePackage{microtype}
\RequirePackage{amsmath,amssymb}
\RequirePackage{xcolor}
\RequirePackage{enumitem}
\RequirePackage[most]{tcolorbox}
\RequirePackage{tabularx,booktabs,array}
\RequirePackage{fancyhdr}
\RequirePackage{needspace}
\RequirePackage{geometry}
\RequirePackage[hidelinks]{hyperref}

\geometry{a4paper,top=17mm,bottom=18mm,left=23mm,right=23mm}
\definecolor{notegray}{gray}{0.39}
\setlength{\parindent}{0pt}
\setlength{\parskip}{2.5pt plus 1pt minus 0.5pt}
\setlist[itemize]{topsep=2pt,itemsep=1pt,parsep=0pt,leftmargin=1.45em}
\setlist[enumerate]{topsep=2pt,itemsep=2pt,parsep=0pt,leftmargin=1.65em}
\setcounter{secnumdepth}{0}
\renewcommand{\arraystretch}{1.12}

% Zahlavi: vlevo rocnik | globalni tema, vpravo cislo: nazev hodiny
% Zapati zustava zcela prazdne.
\pagestyle{fancy}
\fancyhf{}
\renewcommand{\headrulewidth}{0.35pt}
\renewcommand{\footrulewidth}{0pt}
\setlength{\headheight}{14pt}
\newcommand{\nastavPripravu}[4]{%
  \fancyhead[L]{\small #1\enspace\textbar\enspace #2}%
  \fancyhead[R]{\small #3: #4}%
}

% Minimalisticke cernobile boxy. Kazda definice zvlast.
\tcbset{commonlesson/.style={enhanced,breakable=false,colback=white,colframe=black,
  colbacktitle=white,coltitle=black,boxrule=0.45pt,
  arc=0.8mm,outer arc=0.8mm,left=3mm,right=3mm,top=1.6mm,bottom=1.6mm,
  toptitle=1mm,bottomtitle=0.7mm,fonttitle=\bfseries,
  attach boxed title to top left={xshift=2.5mm,yshift=-1.15mm},
  boxed title style={colback=white,colframe=white,boxrule=0pt,left=1mm,right=1mm}}}
\newtcolorbox{zadaniBox}[1][]{commonlesson,title={Zadání},#1}
\newtcolorbox{prikladBox}[1][]{commonlesson,title={Příklad},#1}
\newtcolorbox{definiceBox}[1][]{commonlesson,title={Definice},#1}

% Spolecny box pro zadani a jeho reseni, bez mezery mezi castmi.
% Pouziti:
%   \begin{prikladReseniBox}
%     Zadani ulohy.
%     \tcblower
%     \begin{align*} ... \end{align*}
%   \end{prikladReseniBox}
% Spodni cast ma jemny sedy podklad a delici linku; lze v ni sazet
% vice radku matematiky. Obe casti maji jen jeden spolecny ramecek.
\tcbset{solutionlesson/.style={
  commonlesson,
  enhanced jigsaw,
  breakable,
  colbacklower=black!3,
  segmentation style={solid,draw=black!25,line width=0.35pt},
  fontlower=\small\color{notegray},
  before lower={\noindent\textbf{Řešení:}\enspace},
  leftlower=3mm,rightlower=3mm
}}
\newtcolorbox{prikladReseniBox}[1][]{solutionlesson,title={Příklad},#1}
\newtcolorbox{zadaniReseniBox}[1][]{solutionlesson,title={Zadání},#1}

% Pomucky pro autora pripravy: sedive poznamky; \reseni pro starsi pripravy
\newcommand{\poznamka}[1]{\par\begingroup\small\color{notegray}\itshape
  \textnormal{\textbf{Poznámka pro učitele.}}\enspace #1\par\endgroup}
\newcommand{\reseni}[1]{\par\begingroup\small\color{notegray}
  \textbf{Očekávané řešení:}\enspace #1\par\endgroup}
\newcommand{\blok}[1]{\needspace{5\baselineskip}\subsection*{#1}}
\newcommand{\vyrok}[1]{\uv{#1}}

\usepackage{tikz}
\usetikzlibrary{arrows.meta}
\nastavPripravu{1. ročník}{Funkce}{59}{Co je funkce?}
\begin{document}

\poznamka{\textbf{Jádro na 45 minut:} motivace 4 min, přiřazovací diagramy
7 min, základní definice 6 min, značení 5 min, způsoby zadání 8 min,
samostatné procvičení 10 min a kontrola 5 min. Množinové vyjádření přes
kartézský součin je \textbf{doplňující pohled} pro žáky, nikoli další
povinný výklad v této hodině. Nejde o úvod do lineárních funkcí;
podrobněji budeme obory určovat v hodině 60 a grafy číst v hodině 61.}

\blok{Kdy jedna veličina závisí na druhé?}
Čtverec má stranu délky $a$ a obsah $S=a^2$. Pro stranu délky $2\,\mathrm{cm}$
je obsah $4\,\mathrm{cm}^2$; pro stranu délky $3\,\mathrm{cm}$ je obsah
$9\,\mathrm{cm}^2$. \textbf{Každé přípustné délce strany odpovídá právě jeden
obsah.} Délka strany je \emph{vstup}, obsah je \emph{výstup}.

Podobné přiřazení vzniká při výpočtu ceny podle počtu kusů nebo při určování
vzdálenosti od nuly pomocí absolutní hodnoty. Není přitom důležité, zda
výstupy rostou rovnoměrně, zda se některé opakují či zda pravidlo popisuje
přímku. Rozhodující je \emph{jednoznačnost přiřazení}.

\blok{Co smíme nazvat funkcí?}
Podívej se na dva způsoby přiřazení. V levém diagramu má každý vstup právě
jednu šipku; v pravém má jeden vstup dva různé výstupy.
\begin{center}
\begin{minipage}[t]{0.48\linewidth}\centering
\begin{tikzpicture}[>=Latex,scale=.88,every node/.style={font=\small}]
\draw (0,0) ellipse (.65 and 1.05);
\draw (3.25,0) ellipse (.65 and 1.05);
\node at (0,1.35) {$A$};\node at (3.25,1.35) {$B$};
\foreach \yy/\lab in {.65/-1,0/0,-.65/1} {\fill (0,\yy) circle (2pt);\node[left=5pt] at (0,\yy) {$\lab$};}
\foreach \yy/\lab in {.65/0,-.10/1,-.75/2} {\fill (3.25,\yy) circle (2pt);\node[right=5pt] at (3.25,\yy) {$\lab$};}
\draw[->] (.11,.65)--(3.08,-.10);
\draw[->] (.11,0)--(3.08,.65);
\draw[->] (.11,-.65)--(3.08,-.10);
\end{tikzpicture}

\small\textbf{Funkce.} Různé vstupy mohou mít stejný výstup.
\end{minipage}\hfill
\begin{minipage}[t]{0.48\linewidth}\centering
\begin{tikzpicture}[>=Latex,scale=.88,every node/.style={font=\small}]
\draw (0,0) ellipse (.65 and 1.05);
\draw (3.25,0) ellipse (.65 and 1.05);
\node at (0,1.35) {$A$};\node at (3.25,1.35) {$B$};
\foreach \yy/\lab in {.5/1,-.55/2} {\fill (0,\yy) circle (2pt);\node[left=5pt] at (0,\yy) {$\lab$};}
\foreach \yy/\lab in {.65/a,0/b,-.65/c} {\fill (3.25,\yy) circle (2pt);\node[right=5pt] at (3.25,\yy) {$\lab$};}
\draw[->] (.11,.5)--(3.08,0);
\draw[->] (.11,.5) to[bend left=16] (3.08,.65);
\draw[->] (.11,-.55)--(3.08,-.65);
\end{tikzpicture}

\small\textbf{Není funkce.} Vstup $1$ má dva různé výstupy.
\end{minipage}
\end{center}

\begin{definiceBox}
\textbf{Funkce} $f\colon A\to B$ je přiřazení, které \textbf{každému}
prvku $x\in A$ přiřazuje \textbf{právě jeden} prvek $y\in B$.
Tento prvek zapisujeme $y=f(x)$. Množina $A$ se nazývá
\textbf{definiční obor} funkce a značí se $D(f)$.
\end{definiceBox}

Zápis $f(x)$ čteme \uv{funkční hodnota v bodě $x$}: $f$ označuje celou
funkci, $x$ její \textbf{argument} (vstup) a $f(x)$ odpovídající výstup.
\textbf{Nejde o součin $f\cdot x$.} Výstupy se mohou opakovat, ale jednomu
vstupu nesmějí odpovídat dvě různé hodnoty. Podrobné určování $D(f)$ a
oboru hodnot přijde až v příští hodině.

\blok{Značení funkce}
\begin{center}
\small
\renewcommand{\arraystretch}{1.12}
\begin{tabularx}{\linewidth}{@{}p{28mm}X@{}}
\toprule
\textbf{Zápis} & \textbf{Význam}\\
\midrule
$f$ & Funkce jako celek (celé přiřazení).\\
$f(x)$ & Funkční hodnota v argumentu $x$; nejde o součin $f\cdot x$.\\
$y=f(x)$ & Přiřazení argumentu $x$ odpovídající hodnotě $y$.\\
$f(x)=2x+3$ & Předpis pro výpočet funkční hodnoty.\\
$y=2x+3$ & Tentýž předpis zapsaný pomocí proměnných $x,y$.\\
$f\colon x\mapsto 2x+3$ & Zápis přiřazovacího pravidla.\\
$f\colon A\to B$ & Funkce přiřazuje prvkům množiny $A$ hodnoty z $B$.\\
\bottomrule
\end{tabularx}
\end{center}
Rozlišuj \textbf{funkci} $f$, její \textbf{argument} $x$ a
\textbf{funkční hodnotu} $f(x)$. Samotný vzorec ještě nemusí určit
všechny přípustné argumenty; to budeme řešit v následující hodině.


\blok{Jiný matematický pohled na funkci}
\poznamka{Následující dvě definice ukazují, že funkci lze popsat také jako
množinu uspořádaných dvojic. V běžné 45minutové hodině je stačí stručně
představit; podrobný formální zápis je rozšířením pro zájemce.}

\begin{definiceBox}[title={Definice: kartézský součin}]
\textbf{Kartézský součin} množin $A$ a $B$ je množina všech uspořádaných
dvojic, jejichž první složka patří do $A$ a druhá do $B$:
\[
A\times B=\{(x,y)\mid x\in A,\ y\in B\}.
\]
Na pořadí složek záleží: obecně $(x,y)\neq(y,x)$.
\end{definiceBox}

\begin{definiceBox}[title={Definice: funkce jako množina dvojic}]
\textbf{Funkci} $f$ z množiny $A$ do množiny $B$ lze chápat jako takovou
podmnožinu $A\times B$, pro kterou platí
\[
\forall x\in A\;\exists! y\in B:\ (x,y)\in f.
\]
Každý prvek $A$ tedy vystupuje jako první složka právě jedné dvojice.
Zápis $\exists!$ čteme \uv{existuje právě jeden}. Jde o tutéž
jednoznačnost jako v předchozí definici, jen vyjádřenou formálně.
\end{definiceBox}

Například pro $A=\{1,2\}$ a $B=\{a,b,c\}$ dostaneme
\[
A\times B=\{(1,a),(1,b),(1,c),(2,a),(2,b),(2,c)\}.
\]
Množina $f=\{(1,a),(2,c)\}$ popisuje funkci $A\to B$.
Množina $\{(1,a),(1,b),(2,c)\}$ funkci nepopisuje: prvnímu vstupu
přiřazuje dvě různé hodnoty.

\begin{prikladReseniBox}[title={Příklad 1: jedna funkce, několik argumentů}]
Je dána funkce $f(x)=x^2$ pro $x\in\mathbb R$. Vypočti $f(-3)$,
$f(0)$ a $f(2)$. Mohou mít dva různé argumenty stejnou hodnotu?
\tcblower
$ f(-3)=(-3)^2=9$, $f(0)=0^2=0$, $f(2)=2^2=4$.
Ano: $f(-2)=4=f(2)$. Tím není porušena definice funkce, protože
\emph{každý jednotlivý} argument má stále jen jednu hodnotu.
\end{prikladReseniBox}

\blok{Čtyři způsoby zadání funkce}
Jednu a tutéž funkci můžeme zadat několika způsoby.
Uvažujme $f\colon\mathbb R\to\mathbb R$, kde $f(x)=x^2$.

\begin{enumerate}[label=\textbf{\arabic*.},leftmargin=2em,itemsep=1.5pt]
\item \textbf{Tabulkou:} uvedeme argumenty a odpovídající hodnoty.
U funkce na celém $\mathbb R$ zachycuje konečná tabulka jen vybrané body.
\item \textbf{Předpisem:} $f(x)=x^2$ pro všechna $x\in\mathbb R$.
\item \textbf{Grafem:} znázorníme množinu bodů $(x,f(x))$ v rovině.
\item \textbf{Vlastností:} dvojice $(x,y)$ patří k funkci právě tehdy,
když $x\in\mathbb R$ a $y$ je druhá mocnina čísla $x$; množinově
$f=\{(x,y)\in\mathbb R^2\mid y=x^2\}$.
\end{enumerate}

\begin{center}
\begin{minipage}{.42\linewidth}\centering
\small
\begin{tabular}{@{}c|rrrrr@{}}
$x$ & $-2$ & $-1$ & $0$ & $1$ & $2$\\\midrule
$f(x)$ & $4$ & $1$ & $0$ & $1$ & $4$
\end{tabular}

\vspace{2mm}
\footnotesize Tabulka ukazuje \emph{vybrané} hodnoty,
nikoli všechny reálné argumenty.
\end{minipage}\hfill
\begin{minipage}{.54\linewidth}\centering
\begin{tikzpicture}[x=.78cm,y=.53cm,>=Latex]
\draw[step=1,gray!20,very thin] (-2.5,-.45) grid (2.5,4.5);
\draw[->] (-2.65,0)--(2.75,0) node[right] {\small $x$};
\draw[->] (0,-.5)--(0,4.85) node[above] {\small $y$};
\foreach \i in {-2,-1,1,2} \node[below=2pt,font=\scriptsize] at (\i,0) {$\i$};
\foreach \i in {1,2,3,4} \node[left=2pt,font=\scriptsize] at (0,\i) {$\i$};
\draw[very thick,domain=-2.12:2.12,samples=80] plot (\x,{\x*\x});
\fill (-2,4) circle (2pt);\fill (0,0) circle (2pt);\fill (2,4) circle (2pt);
\node[font=\scriptsize] at (1.9,2.9) {$y=x^2$};
\end{tikzpicture}
\end{minipage}
\end{center}

\poznamka{V této hodině po žácích ještě nechtěj sestrojování paraboly.
Graf je jen jeden ze způsobů zápisu závislosti. V hodině 61
budou podrobně číst souřadnice a rozhodovat, zda obrazec může být grafem funkce.}

\begin{prikladReseniBox}[title={Příklad 2: funkce není vždy druhá mocnina}]
Pro $g(x)=|x|$ a $h(x)=2x+1$ vypočti $g(-3)$ a $h(-3)$.
Rozhodni, zda obě funkce přiřazují argumentu $-3$ stejný výstup.
\tcblower
$g(-3)=|-3|=3$, zatímco $h(-3)=2(-3)+1=-5$.
Obě pravidla jsou funkcemi, ale jsou to \emph{různé funkce}.
Z definice nevyplývá, že různá pravidla musí mít stejné hodnoty.
\end{prikladReseniBox}

\blok{Samostatné procvičení — základ}
\begin{zadaniReseniBox}[title={Samostatná práce}]
\begin{enumerate}[label=\alph*)]
\item Funkce $p$ je na $A=\{1,2,3\}$ zadána dvojicemi
$\{(1;4),(2;4),(3;9)\}$. Je to funkce? Proč?
\item Je přiřazení $\{(1;2),(1;3),(2;4)\}$ funkcí
z $A=\{1,2\}$? Vysvětli.
\item Pro $q(x)=x^2+1$ urči $q(-2)$ a $q(3)$.
\item Pro $r(x)=|x-2|$ urči $r(0)$ a $r(5)$.
\item U funkce $f(x)=x^2$ najdi dva různé argumenty se stejnou hodnotou $9$.
\end{enumerate}
\tcblower
\begin{enumerate}[label=\alph*)]
\item Ano. Každý vstup má právě jeden výstup; opakování hodnoty $4$ nevadí.
\item Ne. Vstupu $1$ odpovídají dvě různé hodnoty, $2$ a $3$.
\item $q(-2)=(-2)^2+1=5$, $q(3)=3^2+1=10$.
\item $r(0)=|0-2|=2$, $r(5)=|5-2|=3$.
\item Například $x=-3$ a $x=3$, protože $(-3)^2=3^2=9$.
\end{enumerate}
\end{zadaniReseniBox}

\poznamka{Volitelná otázka pro rychlé: tabulka s hodnotami $2,1,2$ při
argumentech $-1,0,1$ plně určuje funkci na $\{-1,0,1\}$, ale nikoli na
celém $\mathbb R$. Připravuje to příští hodinu.}

\blok{Domácí procvičení}
\begin{zadaniReseniBox}[title={Domácí úkol}]
\begin{enumerate}[label=\alph*)]
\item Pro $f(x)=x^2-2$ vypočti $f(-3)$, $f(0)$ a $f(3)$.
\item Pro $g(x)=|x+1|$ urči $g(-4)$, $g(-1)$ a $g(2)$.
\item Je $S=\{(-1;1),(0;0),(1;1),(1;2)\}$ grafem funkce na
$\{-1,0,1\}$? Zdůvodni.
\item Funkce $h$ je na $\{0,1,2,3\}$ zadána hodnotami
$h(0)=2$, $h(1)=3$, $h(2)=2$, $h(3)=5$.
Vypiš všechny uspořádané dvojice a najdi argumenty, pro které je $h(x)=2$.
\item Čtverec má stranu $a>0$ v centimetrech a obsah $S(a)=a^2$
v $\mathrm{cm}^2$. Urči $S(5)$ a vysvětli, proč hodnoty $a=-5$ do tohoto modelu nepatří.
\item Pro $u(x)=x^2$, $x\in\mathbb R$, najdi všechny argumenty, kterým
funkce přiřadí hodnotu $16$.
\end{enumerate}
\tcblower
\begin{enumerate}[label=\alph*)]
\item $f(-3)=9-2=7$, $f(0)=-2$, $f(3)=9-2=7$.
\item $g(-4)=|-3|=3$, $g(-1)=0$, $g(2)=|3|=3$.
\item Ne: vstupu $1$ přísluší dvě různé hodnoty, $1$ a $2$.
\item $\{(0;2),(1;3),(2;2),(3;5)\}$; hodnotu $2$ mají argumenty $0$ a $2$.
\item $S(5)=25\,\mathrm{cm}^2$. Délka strany nemůže být záporná;
$a=-5$ není přípustný argument \emph{tohoto modelu}, i když $(-5)^2$ existuje.
\item $x=-4$ nebo $x=4$, protože obě čísla mají druhou mocninu $16$.
\end{enumerate}
\end{zadaniReseniBox}

% Závěrečná kontrola: vysvětlit jednoznačnost, rozlišit f od f(x),
% vyvrátit mylné tvrzení, že různé vstupy musejí mít různé výstupy.
\end{document}
